\documentclass[11pt,a4paper]{article}
\usepackage[utf8]{inputenc}
\usepackage[T1]{fontenc}
\usepackage[spanish]{babel}
\usepackage{amsmath,amssymb,mathtools}
\usepackage{geometry}
\usepackage{enumitem}
\usepackage{lmodern}
\usepackage{microtype}
\geometry{margin=2.4cm}
\setlength{\parindent}{0pt}
\setlength{\parskip}{0.45em}
\newcommand{\R}{\mathbb{R}}
\newcommand{\C}{\mathbb{C}}
\newcommand{\K}{\mathbb{K}}
\newcommand{\Ker}{\operatorname{Ker}}
\newcommand{\rg}{\operatorname{rg}}
\newcommand{\sg}{\operatorname{sg}}
\newcommand{\Span}{\operatorname{L}}
\begin{document}

\begin{center}
{\LARGE\bfseries Álgebra Lineal}\\[0.4em]
{\Large Tema 5: Formas Bilineales y Cuadráticas --- Ejercicios}
\end{center}

\section{Ejercicio 1}
Estudiar si:
\[
f:\R^2\times\R^2\to\R
\]
\[
f((x,y),(x',y'))=xx'+4xy'-3yx'
\]
es una forma bilineal.

\section{Ejercicio 2}
Estudiar si:
\[
f:\R^2\times\R^2\to\R
\]
\[
f((x,y),(x',y'))=2xy'+5x-3y
\]
es una forma bilineal.

\section{Ejercicio 3}
Estudiar si:
\[
f:\R^3\times\R^3\to\R
\]
\[
f((x_1,x_2,x_3),(y_1,y_2,y_3))=x_1y_1+1
\]
es una forma bilineal.

\section{Ejercicio 4}
En el espacio vectorial de funciones continuas en $[a,b]$, estudiar:
\[
f(p,q)=\int_a^bp(x)q(x)\,dx.
\]

\section{Ejercicio 5}
Demostrar que toda forma bilineal puede descomponerse como suma de una forma bilineal simétrica y otra antisimétrica.

\section{Ejercicio 6}
Sea:
\[
f:\R^3\times\R^3\to\R
\]
\[
f(x,y)=x_1y_1-x_1y_3+x_2y_2-x_3y_1+x_3y_2+2x_3y_3+x_2y_3.
\]
Calcular:
\begin{enumerate}[label=\alph*)]
\item La expresión matricial de $f$.
\item El rango y núcleo.
\item Si:
\[
x_1'=2x_1+2x_2+3x_3,\quad
x_2'=-x_1-x_3,\quad
x_3'=-x_2-x_3,
\]
obtener la expresión de $f$ en la base correspondiente.
\end{enumerate}

\section{Ejercicio 7}
De una forma bilineal simétrica:
\[
f:\mathcal P_1(x)\times\mathcal P_1(x)\to\R
\]
se sabe:
\[
f(x+1,x+1)=8,\qquad
f(x+2,x+2)=11,\qquad
f(x,x)=3.
\]
Calcular su matriz respecto de $\{1,x\}$.

\section{Ejercicio 8}
Sea:
\[
f:\mathcal P_2(x)\times\mathcal P_2(x)\to\R,
\qquad
f(p,q)=p(0)q(0).
\]
Obtener $M_B(f)$ para:
\[
B=\{2+x-x^2,\;1+2x,\;3\}.
\]

\section{Ejercicio 9}
Clasificar como simétricas, antisimétricas o ninguna:
\[
\text{a) }f(x,y)=2x_1y_1-5x_1y_2-5x_2y_1+4x_2y_2
\]
\[
\text{b) }g(x,y)=3x_1y_2-3x_2y_1
\]
\[
\text{c) }h(x,y)=x_1y_1-2x_1y_2+7x_2y_1+6x_2y_2.
\]

\section{Ejercicio 10}
Descomponer:
\[
f(x,y)=2x_1y_1+7x_1y_2+6x_2y_1-x_2y_2
\]
como suma de una forma simétrica y otra antisimétrica.

\section{Ejercicio 11}
La matriz de una forma bilineal:
\[
f:E\times F\to\K
\]
en:
\[
B_E=\{u_1,u_2\},
\qquad
B_F=\{v_1,v_2,v_3\}
\]
es:
\[
A=
\begin{pmatrix}
2&-1&1\\
3&4&1
\end{pmatrix}.
\]
Si:
\[
B_E'=\{u_1-u_2,u_1+u_2\}
\]
y:
\[
B_F'=\{v_1,v_1+v_2,v_1+v_2+v_3\},
\]
hallar la matriz de $f$ en estas bases.

\section{Ejercicio 12}
La matriz de:
\[
f:\R^2\times\R^2\to\R
\]
en:
\[
B=\{(1,2),(3,-7)\}
\]
es:
\[
A=
\begin{pmatrix}
2&-1\\
3&6
\end{pmatrix}.
\]
Obtener la matriz en la base canónica.

\section{Ejercicio 13}
Sea $S_2$ el espacio de matrices simétricas $2\times2$, con base:
\[
B=
\left\{
\begin{pmatrix}1&0\\0&0\end{pmatrix},
\begin{pmatrix}0&1\\1&0\end{pmatrix},
\begin{pmatrix}0&0\\0&1\end{pmatrix}
\right\}.
\]
Sea $f$ simétrica con:
\[
M_B(f)=
\begin{pmatrix}
1&1&0\\
1&2&-1\\
0&-1&1
\end{pmatrix}.
\]
\begin{enumerate}[label=\alph*)]
\item Calcular el rango y una base del núcleo.
\item Calcular el subespacio conjugado de:
\[
\begin{pmatrix}1&2\\2&0\end{pmatrix}.
\]
\item Construir una base de vectores conjugados dos a dos.
\item Clasificar la forma cuadrática asociada.
\end{enumerate}

\section{Ejercicio 14}
Sea:
\[
q(x_1,x_2,x_3)
=
x_1^2+5x_2^2+8x_3^2
+4x_1x_2-6x_1x_3-8x_2x_3.
\]
Determinar una base de $\R^3$ respecto de la que la matriz de la forma cuadrática sea diagonal con $1,-1,0$.

\section{Ejercicio 15}
Clasificar según $a$:
\[
q(x_1,x_2,x_3)
=
x_1^2+x_2^2+5x_3^2
+2ax_1x_2-2x_1x_3+4x_2x_3.
\]

\section{Ejercicio 16}
Clasificar según $a$ la forma cuadrática de matriz:
\[
A=
\begin{pmatrix}
1&2&1\\
2&6&2\\
1&2&a
\end{pmatrix}.
\]

\section{Ejercicio 17}
Sea:
\[
q(x_1,x_2,x_3,x_4)
=
x_1^2+3x_2^2-x_3^2
+4x_1x_2+2x_1x_3+2x_1x_4
+2x_2x_4-2x_3x_4.
\]
\begin{enumerate}[label=\alph*)]
\item Diagonalizar y clasificar $q$.
\item Obtener una base del núcleo de la forma polar asociada.
\end{enumerate}

\section{Ejercicio 18}
Diagonalizar y clasificar:
\[
q(x,y,z)=x^2+y^2+3z^2+4xy+2xz+2yz.
\]

\section{Ejercicio 19}
Diagonalizar y clasificar:
\[
q(x,y,z)=x^2-3z^2-2xy+2xz-6yz.
\]

\section{Ejercicio 20}
Diagonalizar y clasificar:
\[
q(x,y,z,t)=xy+yz+xt+yt+zt.
\]

\section{Ejercicio 21}
Diagonalizar y clasificar:
\[
q(x,y,z)=3y^2+2xz.
\]

\section{Ejercicio 22}
Diagonalizar y clasificar según $a$:
\[
q(x,y,z)=x^2+ay^2+3z^2+2xy+2xz+2ayz.
\]

\section{Ejercicio 23}
Diagonalizar y clasificar según $a$:
\[
q(x,y,z)=x^2-ay^2+2xz-4ayz.
\]

\section{Ejercicio 24}
Sea:
\[
q(x,y,z)=x^2+y^2+(k+1)z^2+2xz+2kyz.
\]
\begin{enumerate}[label=\alph*)]
\item Obtener su matriz.
\item Clasificar según $k$.
\item Para $k=2$, obtener el conjugado de:
\[
S=\{(x,y,z):x-y=0,\;y=0\}.
\]
\item Diagonalizar para $k=3$.
\item Para $k=4$, diagonalizar obteniendo una base de vectores conjugados.
\end{enumerate}

\section{Ejercicio 25}
En $\R^3$:
\[
q_a(x_1,x_2,x_3)
=
x_1^2+4x_2^2+2x_3^2+2x_1x_2+4ax_1x_3.
\]
\begin{enumerate}[label=\alph*)]
\item Obtener la matriz asociada en la base canónica.
\item Determinar si existen $\lambda,\mu\in\R$ para que:
\[
B=
\{u_1=e_1,\;
u_2=e_1+\lambda e_2,\;
u_3=ae_1+ae_2+\mu e_3\},
\qquad a\neq0,
\]
sea una base de vectores conjugados.
\item Determinar para qué valores de $a$ es definida positiva o semidefinida positiva.
\item Calcular:
\[
\dim\left[\Span\{u_3\}\right]^c
\]
para $a=1$.
\item Para $a=2$, diagonalizar por dos métodos distintos, indicando $P$, $D$ y la signatura.
\end{enumerate}

\end{document}
